\scanpage{31}
Supposons qu'il y a des territoires
constitués de blanc et noir, considérés
\unsure{indélogeables}, blanc a joué dernier
(alors c'est à noir de jouer) on a
joué $n$ coups
\[
n=p+q\qquad\deleted{\unsure{n}}n'=p'+q'\qquad n=n'
\]
\margin{$p$ : nb de pions noirs sur goban
qui ne sont prisonniers virtuels}

$q$ nb de prisonniers virtuels

$t$ territoire noir (\unsure{avec pris enlevé} les pris. bl.)

$t'$ \ldots\ blanc (\text{---} noirs)
\margin{$p+p'+t+t'=\unsure{v|\nu}^2$}
\[
\begin{array}{ll}
\deleted{\unsure{Q}}\ s=t+q'&\text{pts de noir potentiels}\\
s'=t'+q&\text{---}\quad\text{blanc}\quad\text{---}
\end{array}
\]
\[
d=s-s'=(t-t'\unsure{,}+(q'-q)=(t-t')+(p-p')
\]
\[
\unsure{v|\nu}^2+d=2(p+t)\qquad\text{i.e.}\qquad\boxed{d=2(p+t)-\unsure{v|\nu}^2}
\]
\[
D=d\unsure{-|+}\varepsilon\qquad\varepsilon=0,1\qquad(\mathrm{NB}\quad p+t=P+T)
\]
\margin{?}
\figure{\begin{tikzpicture}\draw (0,1.8)--(0,.4)..controls (0,.15) and (.15,0)..(.3,0);\end{tikzpicture}}
Comme on ne peut prévoir la valeur de
$\varepsilon$ (est-elle indépendante de la suite du
jeu?) il est raisonnable de compenser
l'avantage de noir (de commencer)
en posant $\varepsilon=1$ (i.e. en admettant que
noir joue dernier)
\mspage{30}
